\documentclass[11pt,a4paper]{article}
\usepackage[margin=2.3cm]{geometry}
\usepackage{amsmath,amssymb}
\usepackage{booktabs}
\usepackage{array}
\usepackage{graphicx}
\usepackage{microtype}
\usepackage{tabularx}
\usepackage{xcolor}
\usepackage{enumitem}
\usepackage[hidelinks]{hyperref}
\pdfstringdefDisableCommands{\def\code#1{#1}}
\setlist{itemsep=2pt,topsep=4pt}
\newcommand{\Lag}{\mathscr{L}}
\usepackage{mathrsfs}
\newcommand{\code}[1]{\texttt{#1}}
\newcommand{\winf}{w_\infty}
\renewcommand{\arraystretch}{1.15}

\title{DG--Laguerre schemes for the Euler equations on a half-space:\\[2pt]
\large primitive vs.\ conservative variables, and two treatments of the unbounded direction}
\author{}

\emergencystretch=3em
\begin{document}
\maketitle

\begin{abstract}
\noindent This report compares two design choices for matrix-free DG--Laguerre
discretizations of the Euler equations with reaction on
$\Omega=\Omega_b\times[0,\infty)$: (i) primitive variables $(\rho,u,E)$ versus
conservative variables $(\rho,\rho u,E)$ as expanded unknowns, and (ii) the
treatment of the unbounded direction $z$ in \code{laguerre\_systems.cc} (code A, following the notes
\emph{laguerre\_systems.pdf}: pure (standalone)
Laguerre basis, exact $z$-integration) versus \code{laguerre\_systems\_2.cc} (code B, following the notes
\emph{laguerre\_systems\_2.pdf}: DG--Legendre
intervals plus a Laguerre tail, numerical (inexact) quadrature). In A we project the $d$-system onto a $(d-1)-$system by integrating in $(0,\infty)$ first; in this case we compute the Rusanov flux of the projected system, with one $\lambda$ per $x$-face point. In B the Rusanov flux is evaluated pointwise at every $z$-point with local $\lambda(z)$, and \emph{then} integrated. All numbers come from runs
of the two programs on the same machine, single core, 2D ($d=2$, $\Omega_b=(0,10)$,
80 cells, DG $Q_3$ in $x$) unless stated otherwise.

\smallskip\noindent\textbf{Main findings.}
\begin{itemize}
\item The accuracy gap between the codes with pure Laguerre bases comes from the
  \emph{choice of variables}, not from exact integration. In primitive variables the
  quadrature-based code reproduces the errors of the exact-integration code to
  within 4\%.
\item In a pure Laguerre basis, primitive variables are 1.5--2$\times$ more
  accurate on the moving vortex. The reason is approximation: the momentum and
  energy perturbations contain narrower structures than the velocity perturbation.
  With a local (DG-in-$z$) basis this advantage disappears.
\item Primitive variables cost an extra mass-matrix solve: $2.6$--$3.9\times$ the run
  time with pure Laguerre and $1.6$--$1.9\times$ with the hybrid basis. Stability limits are
  identical.
\item For the $z$-direction, the hybrid DG--Legendre $+$ Laguerre-tail basis of
  code B is the better choice. It reaches $7.7\cdot10^{-6}$ in
  3.9\,s, against $2.6\cdot10^{-5}$ in 25.8\,s for the pure Laguerre code. Exact
  $z$-integration buys no measurable accuracy over pointwise quadrature, and each
  time step costs about $5\times$ more. (i.e. the Propositions with quadrature results listed in \emph{laguerre\_systems.pdf} are not that useful in practice, approximate quadrature is almost as accurate and does not require storage of several matrices)
\end{itemize}
\end{abstract}

\section{Setting, codes and test problems}

\paragraph{Equations.} We solve
$\partial_t(M(w)w)+\nabla\cdot F(w)=b(w,x,z,t)$ with $w=(\rho,u,E)$,
$M(w)w=(\rho,\rho u,E)$, $p=(\gamma-1)(E-\tfrac12\rho|u|^2)$, $\gamma=1.4$, and
$b=(0,\rho g,\rho u\cdot g)-\Gamma(z)\,(M(w)w-M(\winf)\winf)$, where
$\Gamma_{ii}(z)=\Delta\gamma_i/(1+e^{(\alpha_i-z)/\sigma_i})$ is a sigmoidal
damping layer. Solutions tend to the far-field state $\winf$ as $z\to\infty$. The
perturbation from the far field is expanded in $z$,
\begin{equation}\label{eq:exp}
V(x,z,t)\approx V_\infty+\sum_j W_j(x,t)\,\psi_j(z),
\end{equation}
where $V$ is either the primitive vector $w$ or the conservative vector $U=M(w)w$.
Each $W_j$ is a DG function on $\Omega_b$. Testing with $\psi_i$ and integrating
over $z$ gives the projected system
$\partial_t\tilde a+\nabla_x\cdot\tilde B=\tilde h$ of the notes, which is
discretized with matrix-free DG in $x$ and explicit low-storage RK(5,4) in time.

\filbreak
\paragraph{Implementations compared.}
\begin{center}
\begin{tabularx}{\textwidth}{@{}lX@{}}
\toprule
\textbf{A} & \code{laguerre\_systems.cc}: primitive variables, pure Laguerre basis
  $\psi_j=\Lag^\mu_j=e^{-\mu z/2}L_j(\mu z)$, exact $z$-integrals (Proposition 2),
  Rusanov flux of the projected system, momentum mass solved by preconditioned CG.\\
\textbf{B$_c$} & \code{laguerre\_systems\_2.cc} with conservative variables (default):
  DG--Legendre on $N_z$ intervals of $[0,L_z]$ plus a Laguerre tail on $[L_z,\infty)$
  ($N_z=0$: pure Laguerre), pointwise $z$-quadrature, local Lax--Friedrichs flux at
  every $z$-point, diagonal mass.\\
\textbf{B$_p$} & \code{laguerre\_systems\_2.cc} with primitive variables (option
  \code{primitive}): as B$_c$, with the momentum mass solved directly
  (Section~\ref{sec:primitive-impl}).\\
\bottomrule
\end{tabularx}
\end{center}

\paragraph{Test problems.}
\begin{itemize}
\item \emph{Vortex along $x$}: the isentropic vortex of step-67 in the $(x,z)$ plane,
  centered at $(3,2)$ and advected with far-field velocity $u_\infty=e_x$. It is an
  exact solution without source, with exact Dirichlet data on all boundaries. We
  report the $L^2(\Omega_b\times(0,\infty))$ error of the density at $t=1$, the same
  metric in all codes, computed with a fine $z$-rule.
\item \emph{Vortex along $+z$}: center $(5,2)$, $u_\infty=e_z$, up to $t=6$. The
  vortex rises into a damping layer ($\Delta\gamma=3$, $\alpha=4.5$, $\sigma=0.4$,
  all components) or, without damping, into the region the $z$-basis cannot
  resolve. This is a qualitative and long-time stability test.
\item \emph{Manufactured solution} (B only): smooth perturbation decaying like
  $e^{-z}$, constant reaction $0.5$, gravity $g=-0.5\,e_z$, $t=0.1$, mesh
  $(0,1)^{d-1}$.
\end{itemize}

%=====================================================================
\section{Primitive vs.\ conservative variables}
%=====================================================================

\subsection{Theory}

\paragraph{1. Structure of the semi-discrete system.}
Both formulations discretize the same balance law in conservation form; only the
expanded variable $V$ in \eqref{eq:exp} differs.
\begin{itemize}
\item \emph{Conservative} ($V=U$): $\tilde a$ is linear in the unknowns, so the
  semi-discrete system is $\mathbb M\,\dot W=R(W)$ with the constant matrix
  $\mathbb M=M_x\otimes M_z$. Here $M_z$ is diagonal, because the $\psi_j$ are
  orthogonal ($\int\Lag_i^\mu\Lag_j^\mu=\delta_{ij}/\mu$, and orthonormal Legendre
  polynomials on each interval). An explicit stage costs one cell-wise inverse mass
  application.
\item \emph{Primitive} ($V=w$): $\tilde a$ is nonlinear in $W$. The density and
  energy rows are unchanged, $\partial_t\int\rho\,\psi_i=m_i\dot\rho_i$, but the
  momentum rows become
  \begin{equation}\label{eq:mom}
  \sum_k\Big[\int\rho\,\psi_i\psi_k\,dz\Big]\dot u_k
  \;=\;R^m_i-\int\dot\rho\,u\,\psi_i\,dz ,
  \end{equation}
  where $\dot\rho$ is known from the density equation. The matrix in \eqref{eq:mom}
  is a $\rho$-weighted mass matrix. It is SPD if and only if $\rho>0$, depends on the
  state, and couples all $z$-modes, and through $x$ all DoFs of a cell. Every RK
  stage therefore needs a cell-local linear solve.
\end{itemize}

\paragraph{2. Polynomial structure of the fluxes.}
In primitive variables $M(w)w$, $F(w)$ and the gravity source are polynomials of
degree at most $4$ in the entries of $w$; for example $(E+p)u=\gamma Eu-\tfrac{\gamma-1}{2}\rho|u|^2u$.
This is what makes Proposition 2 of \emph{laguerre\_systems.pdf} possible: products of Laguerre
expansions are integrated \emph{exactly} by Gauss--Laguerre rules with scaling
$(s+1)\mu/2$, $s=0,\dots,4$. In conservative variables $u=m/\rho$ and $p$ are
rational functions of $U$, so no finite rule is exact and the $z$-integrals always
carry a quadrature (aliasing) error. Section~\ref{sec:num-vars} shows that, at the
resolutions used here, this error is \emph{not} the dominant one.

\paragraph{3. Conservation.}
Both formulations are conservative in $x$: each Laguerre moment satisfies a
flux-form equation, and the single-valued numerical fluxes telescope across
interior faces. In the primitive formulation the conserved quantity is the discrete
$\tilde a(W)$, evaluated with the same quadrature as the mass term. In B$_p$ the
$\rho$-weight is collocated at the $p+1$ Gauss points in $x$. Since the DG basis is a
partition of unity, summing the test functions shows that
$\sum_q J_q w_q\,\partial_t(\rho u)(x_q)$ changes only through boundary fluxes. In
$z$, exact global conservation would require the constant function in the test
space. The Laguerre functions do not contain it, whereas the DG intervals of B do
contain it locally. This limitation is common to both choices of variables.

\paragraph{4. Approximation properties.}
The discrete solution lives in $V_\infty+\operatorname{span}\{\psi_j\}$, so accuracy
depends on how well the basis approximates the \emph{chosen} variables. For the
moving vortex let $f\propto e^{-|x-x_c|^2}$, a Gaussian of width $1$. Then
\begin{gather*}
u'=u-u_\infty\propto f\cdot(z-z_c),\qquad
\rho'=(1-c f^2)^{1/(\gamma-1)}-1\sim f^2,\\
m'=\rho u-\rho_\infty u_\infty=\rho_\infty u'+u_\infty\rho'+\rho'u'.
\end{gather*}
Because the far-field velocity is nonzero, the momentum perturbation contains
$u_\infty\rho'$, which has the narrower width $1/\sqrt2$ of $f^2$. The energy
perturbation contains $\rho|u|^2$ and is similarly narrow. The velocity perturbation
has the width of $f$.

A global Laguerre basis converges spectrally, with a rate that depends strongly on
the width and analyticity of the function. Its resolution also degrades with height,
since the node spacing grows like $\sqrt{z/(\mu p_z)}$. Narrower functions are
therefore approximated much worse. A local piecewise-polynomial basis converges
algebraically, and the width enters only through a constant.

A related remark follows from Remark 6 of the notes. If each perturbation decays
like $e^{-\mu z/2}$, then the products in $m'$ and $E'$ decay like $e^{-\mu z}$,
faster than the basis. In conservative variables these products are part of the
expanded unknown and must be approximated in the basis. In primitive variables they
appear only inside the flux integrals, where Proposition 2 handles them exactly.

\paragraph{5. Cost.}
\begin{itemize}
\item Conservative: a diagonal-in-$z$ inverse mass, $O(\#\text{DoFs})$ per stage.
\item Primitive, B$_p$\label{sec:primitive-impl}: the $\rho$-weight is collocated at
  the $p+1$ Gauss points in $x$, where the cell mass matrix is exact. System
  \eqref{eq:mom} then splits into one dense SPD system of size $n_s$, the number of
  modes of segment $s$, per $x$-point and segment. In $z$, the $\rho_\infty$ part is
  the exact diagonal mass and the $\rho-\rho_\infty$ part uses the segment's
  $z$-quadrature. Assembly costs $O(Q_s n_s^2)$ and the Cholesky factorization
  $O(n_s^3/3)$; both are dominated by the Laguerre tail.
\item Primitive, A: the $\rho$-weighted mass uses Gauss($p+3$) in $x$ and is solved by
  CG, preconditioned by the same collocated operator. It needs 3--4 iterations.
\end{itemize}

\paragraph{6. Robustness.}
The primitive formulation requires $\rho>0$ for an SPD mass matrix. Both codes bound
the Cholesky pivot from below. Both formulations need $\rho,p>0$ for the wave speeds
in the fluxes and in the time-step estimate.

\subsection{Numerical comparison}\label{sec:num-vars}

\paragraph{Manufactured solution (B, hybrid basis, $L_z=2$, $p_z=20$).}
Free-stream preservation holds to $|R(0)|=5.9\cdot10^{-13}$ in both formulations.
Errors are $L^2$ errors of all conservative components.
\begin{center}
\begin{tabular}{@{}rrcccc@{}}
\toprule
cells$_x$ & $N_z$ & B$_c$ error & rate & B$_p$ error & rate\\
\midrule
2 & 4 & $4.34\cdot10^{-3}$ & -- & $4.48\cdot10^{-3}$ & --\\
4 & 8 & $3.48\cdot10^{-4}$ & 3.64 & $4.51\cdot10^{-4}$ & 3.31\\
8 & 16 & $2.11\cdot10^{-5}$ & 4.05 & $2.33\cdot10^{-5}$ & 4.27\\
16 & 32 & $1.33\cdot10^{-6}$ & 3.99 & $1.43\cdot10^{-6}$ & 4.03\\
\bottomrule
\end{tabular}
\end{center}
Both formulations converge at the optimal rate 4 for $Q_3$. Here the primitive
errors are 3--30\% larger, and its run time is about 10\% higher. 3D ($Q_2$) gives the
same picture: $3.09\cdot10^{-2}\to4.70\cdot10^{-3}$ conservative and
$3.15\cdot10^{-2}\to4.62\cdot10^{-3}$ primitive.

\paragraph{Vortex along $x$: accuracy and cost.}
Density error at $t=1$; wall times in parentheses. B is run at Courant number 1.5.
\begin{center}
\resizebox{\textwidth}{!}{\begin{tabular}{@{}lccc@{}}
\toprule
$z$-basis & B$_c$ (conservative) & B$_p$ (primitive) & A (primitive, exact)\\
\midrule
Laguerre $p_z=24$, $\mu=10$ & $1.49\cdot10^{-3}$ (2.6\,s) & $8.57\cdot10^{-4}$ (6.7\,s) & $8.51\cdot10^{-4}$ (5.6\,s)\\
Laguerre $p_z=32$, $\mu=10$ & $5.13\cdot10^{-4}$ (4.8\,s) & $2.41\cdot10^{-4}$ (14.9\,s) & $2.33\cdot10^{-4}$ (10.1\,s)\\
Laguerre $p_z=48$, $\mu=10$ & $4.00\cdot10^{-5}$ (12.6\,s) & $2.62\cdot10^{-5}$ (49.1\,s) & $2.61\cdot10^{-5}$ (25.8\,s)\\
\midrule
hybrid $N_z=8$, $L_z=6$, tail $p_z=16$ & $2.93\cdot10^{-3}$ (1.3\,s) & $6.28\cdot10^{-3}$ (2.5\,s) & --\\
hybrid $N_z=16$ & $3.19\cdot10^{-4}$ (2.1\,s) & $4.61\cdot10^{-4}$ (3.5\,s) & --\\
hybrid $N_z=32$ & $7.71\cdot10^{-6}$ (3.9\,s) & $9.19\cdot10^{-6}$ (6.2\,s) & --\\
\bottomrule
\end{tabular}}
\end{center}
Two observations stand out.
\begin{itemize}
\item With pure Laguerre, primitive variables are $1.5$--$2.1\times$ more accurate.
  B$_p$ and A agree to within 4\% although B$_p$ uses inexact pointwise quadrature
  and A exact integration.
\item With the hybrid basis the ranking reverses: conservative variables are
  $1.2$--$2.1\times$ more accurate.
\end{itemize}

\paragraph{Is it quadrature aliasing?}
We repeated the runs with more $z$-points: Laguerre tail $60\to150$ points
($p_z=24$), $76\to150$ ($p_z=32$), and interval rule $5\to9$ points (hybrid
$N_z=16$).
\begin{center}
\begin{tabular}{@{}lcc@{}}
\toprule
& B$_c$: default $\to$ refined & B$_p$: default $\to$ refined\\
\midrule
Laguerre $p_z=24$ & $1.49\cdot10^{-3}\to1.58\cdot10^{-3}$ & $8.57\cdot10^{-4}\to8.45\cdot10^{-4}$\\
Laguerre $p_z=32$ & $5.13\cdot10^{-4}\to4.83\cdot10^{-4}$ & $2.41\cdot10^{-4}\to2.31\cdot10^{-4}$\\
hybrid $N_z=16$ & $3.19\cdot10^{-4}\to3.19\cdot10^{-4}$ & $4.61\cdot10^{-4}\to4.61\cdot10^{-4}$\\
\bottomrule
\end{tabular}
\end{center}
The errors change by at most 6\%. Aliasing of the rational conservative fluxes is
therefore not the cause of the gap. This is consistent with A $\approx$ B$_p$:
exactness of the $z$-integrals does not matter at these resolutions.

\paragraph{Best-approximation errors.}
To test the approximation explanation of item 4, we computed the $L^2$ projection
error of the initial perturbation of each variable onto the $z$-basis, integrated
over $\Omega_b\times(0,40)$.
\begin{center}
\resizebox{\textwidth}{!}{\begin{tabular}{@{}lcccccc@{}}
\toprule
$z$-basis & $\rho'$ & $u'_x$ & $u'_z$ & $E'$ & $(\rho u_x)'$ & $(\rho u_z)'$\\
\midrule
Laguerre $p_z=24$ & $2.8\cdot10^{-4}$ & $1.9\cdot10^{-5}$ & $1.8\cdot10^{-6}$ & $2.4\cdot10^{-3}$ & $1.1\cdot10^{-3}$ & $3.0\cdot10^{-4}$\\
Laguerre $p_z=32$ & $3.6\cdot10^{-5}$ & $3.6\cdot10^{-7}$ & $3.6\cdot10^{-8}$ & $4.3\cdot10^{-4}$ & $2.1\cdot10^{-4}$ & $4.9\cdot10^{-5}$\\
Laguerre $p_z=48$ & $1.7\cdot10^{-6}$ & $1.4\cdot10^{-10}$ & $2.4\cdot10^{-11}$ & $4.1\cdot10^{-5}$ & $1.3\cdot10^{-5}$ & $1.9\cdot10^{-6}$\\
hybrid $N_z=16$ & $7.8\cdot10^{-5}$ & $1.6\cdot10^{-4}$ & $5.4\cdot10^{-5}$ & $5.6\cdot10^{-4}$ & $2.7\cdot10^{-4}$ & $8.0\cdot10^{-5}$\\
hybrid $N_z=32$ & $4.8\cdot10^{-6}$ & $1.0\cdot10^{-5}$ & $3.4\cdot10^{-6}$ & $3.4\cdot10^{-5}$ & $1.6\cdot10^{-5}$ & $5.2\cdot10^{-6}$\\
\bottomrule
\end{tabular}}
\end{center}
In the Laguerre basis the velocity is approximated 50 to $10^5$ times better than
the momentum, as predicted. The momentum error is of the order of the density
error, because $m'\supset u_\infty\rho'$. In the hybrid basis the velocity is even
slightly \emph{worse} than the momentum: $u'$ has amplitude $\approx1$ against
$\approx0.5$ for $\rho'$, and with local polynomials the width only changes a
constant. The simulation errors follow the same ordering. The remaining hybrid
deficit of B$_p$ is plausibly also influenced by the collocated (inexact)
$x$-integration of the $\rho$-weight, which we did not isolate.

\paragraph{Stability.}
The stability limit of the Courant number (vortex along $x$, $t=1$, B with the
spectral-norm time-step estimate) is the same for both variable sets: between 2.5
and 3 with the hybrid basis, and above 4 with pure Laguerre $p_z=48$. In the
long-time vortex-along-$z$ runs ($t=6$), all eight configurations are stable in both
formulations. These include two deliberately under-resolved ones (pure Laguerre
$p_z=16$; hybrid $L_z=3$, $N_z=6$, tail $p_z=8$) in which the vortex rises into
unresolved regions.

\paragraph{Long-time damped vortex (B, $t=6$).}
$\|\rho-\rho_\infty\|$ over $\Omega$, and the deviation from the undamped exact
solution below the layer, $z<3$.
\begin{center}
\resizebox{\textwidth}{!}{\begin{tabular}{@{}lcccc@{}}
\toprule
& \multicolumn{2}{c}{$\|\rho-\rho_\infty\|$, damped, $t=3$ / $t=6$} &
  \multicolumn{2}{c}{$z<3$ deviation, undamped, $t=6$}\\
& B$_c$ & B$_p$ & B$_c$ & B$_p$\\
\midrule
Laguerre $p_z=48$ & 0.1744 / 0.03188 & 0.1744 / 0.03188 & $4.28\cdot10^{-3}$ & $4.24\cdot10^{-3}$\\
hybrid $N_z=24$, tail $p_z=16$ & 0.1744 / 0.03187 & 0.1744 / 0.03187 & $2.55\cdot10^{-5}$ & $2.58\cdot10^{-5}$\\
under-res.\ Laguerre $p_z=16$ & 0.1796 / 0.03580 & 0.1775 / 0.03663 & $1.41\cdot10^{-1}$ & $1.78\cdot10^{-1}$\\
under-res.\ hybrid & 0.1755 / 0.03198 & 0.1751 / 0.03232 & $3.08\cdot10^{-2}$ & $2.80\cdot10^{-2}$\\
\bottomrule
\end{tabular}}
\end{center}
At resolved settings the damping results coincide to 3--4 digits. Run times for the
undamped plus damped pair: Laguerre $p_z=48$ 179\,s (B$_c$) against 669\,s (B$_p$);
hybrid 28.6\,s against 45.7\,s.

\subsection{Conclusions on the variables}
\begin{itemize}
\item \textbf{Accuracy} is decided by approximation, not by integration. With a pure
  Laguerre basis and nonzero far-field velocity, primitive variables are clearly
  better: $1.5$--$2.1\times$ smaller error at equal $p_z$. With a local DG basis in $z$, conservative
  variables are as good or better.
\item \textbf{Cost}: primitive variables need a state-dependent mass solve every
  stage, which makes them $2.6$--$3.9\times$ slower in the pure Laguerre case. At
  equal accuracy B$_p$ with $p_z=32$ ($2.4\cdot10^{-4}$, 14.9\,s) is still slower than
  B$_c$ with $p_z=48$ ($4.0\cdot10^{-5}$, 12.6\,s), which is more accurate.
\item \textbf{Stability and robustness}: no difference was observed. Primitive
  variables additionally require $\rho>0$ for the mass matrix.
\item \textbf{Recommendation}: conservative variables as the default, especially
  with the hybrid basis. Primitive variables are worthwhile only for a pure Laguerre
  discretization where the number of modes, and hence memory, is the bottleneck, or
  where exact integration (Proposition 2) is required.
\end{itemize}

%=====================================================================
\section[Treatment of the z-direction: A vs. B]{Treatment of the $z$-direction:\\ \code{laguerre\_systems.cc} (A) vs.\ \code{laguerre\_systems\_2.cc} (B)}
%=====================================================================

\subsection{Theory}

\begin{center}
\small
\begin{tabularx}{\textwidth}{@{}>{\raggedright\arraybackslash}p{2.6cm}>{\raggedright\arraybackslash}X>{\raggedright\arraybackslash}X@{}}
\toprule
& \textbf{A: \code{laguerre\_systems.cc}} & \textbf{B: \code{laguerre\_systems\_2.cc}}\\
\midrule
$z$-basis & Global Laguerre functions $\Lag^\mu_j$, $j=0..p_z$, on $[0,\infty)$; no faces in $z$. &
  $N_z$ intervals of $[0,L_z]$ with orthonormal Legendre polynomials of degree $p^b_z$
  (DG in $z$), plus a tail $e^{-y}L_i(2y)$, $y=\mu(z-L_z)/2$; pure Laguerre if $N_z=0$.\\
Resolution in $z$ & Node spacing $\sim\sqrt{z/(\mu p_z)}$, degrading with height;
  one parameter pair $(\mu,p_z)$ for the whole column. &
  Uniform $h/p$-refinable resolution on $[0,L_z]$; the tail only has to carry the
  smooth decay to the far field.\\
$z$-integrals & Exact (Proposition 2): polynomial fluxes in $w$ evaluated on graded
  numbers; one Gauss--Laguerre rule per number $s=1..4$ of perturbation factors with
  scaling $(s+1)\mu/2$ and $Q_s=\lceil(1+(s+1)p_z)/2\rceil$ points. $\int f\Lag_i'$
  via the recurrence $\Lag_i'=-\frac\mu2\Lag_i-\mu\sum_{l<i}\Lag_l$. &
  Pointwise quadrature: Gauss($p^b_z+2$) per interval, Gauss--Laguerre with
  $2(p_z+1)+10$ points in the tail. Never exact for nonlinear fluxes. In the tail even
  $\int\psi_i\psi_j$ is inexact: the integrand decays like $e^{-2y}$ but the rule
  weight is $e^{-y}$ (Remark 6 of the notes). The analytic mass $1/\mu$ is used.\\
Coupling in $z$ & Implicit in the global basis; $\partial_z$ through the recurrence. &
  Upwind-type LLF fluxes at interval interfaces and at $z=L_z$ (column-local cell
  kernel).\\
$z=0$ boundary & Pointwise LLF with $\Lag_i(0)=1$; Dirichlet or slip wall. &
  Same, on the trace of the first segment.\\
Lateral faces ($x$) & Rusanov flux of the projected system: one $\lambda$ per
  $x$-face point, the maximum over both sides and all $z$-nodes; exact integrals of
  $F\cdot n$ and $\tilde a$. &
  LLF evaluated pointwise at every $z$-point with local $\lambda(z)$, then
  integrated. Less dissipative where wave speeds are small.\\
Initial / boundary data & Projection with the rule of the notes in $y=\mu z/2$; the
  notes' $Q=p_z+1$ limits the error to $\approx2\cdot10^{-3}$, so $p_z+41$ points are
  used. Lateral Dirichlet data are projected onto the modes. &
  $L^2$ projection with a finer rule; lateral Dirichlet data are evaluated
  pointwise at the $z$-points.\\
Source term & Gravity exact; $\Gamma(z)$ through precomputed matrices
  $\int\Gamma\Lag_j\Lag_k$ plus a $\Gamma$-weighted $s=2$ rule for
  $\rho'u'$; specialized to $\Gamma=\Gamma(z)$. &
  Pointwise $b(U,x,z,t)$ at the quadrature points; any $\gamma(x,z,t)$ and $g$.\\
Mass / time derivative & Primitive variables only; CG with collocated preconditioner. &
  Conservative (diagonal) or primitive (direct collocated solve).\\
Time step & $\mu\|C_z\|_2$ with $C_z(k,j)=\int\Lag_j\Lag_k'$; recomputed every 10 steps. &
  Spectral norm of the 1D upwind DG/Laguerre operators per segment; recomputed every
  10 steps.\\
Code size & 2804 lines & 1895 lines\\
\bottomrule
\end{tabularx}
\end{center}

\paragraph{Consequences.}
\begin{itemize}
\item \emph{Accuracy.} With exact integration the scheme is a true Galerkin
  projection in $z$: no aliasing, and the $z$-discretization of linearized
  operators inherits the structure of the continuous one. B commits a quadrature
  error. It converges spectrally in the number of points but, for rational fluxes,
  never vanishes. The basis matters more than the integration: a global Laguerre
  basis must resolve every structure at the height where it occurs, with a single
  scaling $\mu$. The hybrid basis places resolution where it is needed and leaves the
  far-field decay to the tail.
\item \emph{Cost per step.} A evaluates the fluxes on about
  $\sum_s Q_s\approx\frac72p_z$ nodes ($340$ at $p_z=48$), in truncated polynomial
  arithmetic, on cells and twice on faces. It also needs the CG solve. B evaluates
  each flux once per $z$-point (108 points at $p_z=48$).
\item \emph{Stability.} A has no $z$-aliasing and a strongly dissipative lateral flux.
  B relies on over-integration plus local LLF. In practice both were stable in all
  tests, including under-resolved ones (see below).
\item \emph{Flexibility.} A is tailored to the Euler equations and a $z$-dependent
  sigmoidal $\Gamma$. Exactness requires polynomial fluxes and the specialized
  matrices. B accepts any source and reaction $\gamma(x,z,t)$, and either set of
  variables.
\end{itemize}

\subsection{Numerical comparison}

\paragraph{Exactness checks of A.}
The Proposition-2 rules reproduce $\mu\int\Lag_i\Lag_j=\delta_{ij}$ to
$10^{-13}$. Enlarging any rule by 8 points changes the products of $s$ functions by at
most $2.5\cdot10^{-13}$, for $p_z\le48$.

The $z$-integrals of the source were compared with a brute-force composite
Gauss--Legendre quadrature. Gravity is exact to $10^{-14}$. The reaction integrals,
which cannot be exact because $\Gamma$ is not a polynomial, have relative errors
$5.4\cdot10^{-6}$, $2.3\cdot10^{-6}$, $8.8\cdot10^{-8}$ and $3.0\cdot10^{-9}$ for
$p_z=8,16,32,48$.

\paragraph{Inexactness of B in the tail.}
With the default rule of $2(p_z+1)+10$ points, the computed tail mass matrix
$\int\psi_i\psi_j$ deviates from $\delta_{ij}/\mu$ by a relative
$2.4\cdot10^{-6}$, $1.5\cdot10^{-5}$ and $6.7\cdot10^{-5}$ for $p_z=12,20,32$. The
deviation grows with $p_z$ at fixed over-integration.

\paragraph{Same problem, same metric (vortex along $x$, density error, $t=1$).}
\begin{center}
\resizebox{\textwidth}{!}{\begin{tabular}{@{}lccc@{}}
\toprule
& error & time & steps\\
\midrule
A, Laguerre $p_z=24$ / 32 / 48 & $8.51\cdot10^{-4}$ / $2.33\cdot10^{-4}$ / $2.61\cdot10^{-5}$ & 5.6 / 10.1 / 25.8\,s & 160 / 184 / 233\\
B$_p$, Laguerre $p_z=24$ / 32 / 48 & $8.57\cdot10^{-4}$ / $2.41\cdot10^{-4}$ / $2.62\cdot10^{-5}$ & 6.7 / 14.9 / 49.1\,s & 358 / 424 / 560\\
B$_c$, Laguerre $p_z=24$ / 32 / 48 & $1.49\cdot10^{-3}$ / $5.13\cdot10^{-4}$ / $4.00\cdot10^{-5}$ & 2.6 / 4.8 / 12.6\,s & 358 / 424 / 560\\
B$_c$, hybrid $N_z=8$ / 16 / 32 & $2.93\cdot10^{-3}$ / $3.19\cdot10^{-4}$ / $7.71\cdot10^{-6}$ & 1.3 / 2.1 / 3.9\,s & 194 / 208 / 239\\
\bottomrule
\end{tabular}}
\end{center}
This table separates the three effects.
\begin{enumerate}
\item \emph{Exact vs.\ pointwise $z$-integration} (A vs.\ B$_p$: same variables,
  same basis): no measurable accuracy difference, at most 4\%.
\item \emph{Variables} (B$_p$ vs.\ B$_c$): a factor 1.5--2.1 in favor of primitive
  variables with pure Laguerre (Section~2).
\item \emph{Basis} (Laguerre vs.\ hybrid, B$_c$): the hybrid basis with $N_z=32$ is
  $3.4\times$ more accurate than A's best run and $6.6\times$ faster.
\end{enumerate}

\paragraph{Cost per step} (80 cells, $p_z=48$, pure Laguerre): A 0.111\,s,
B$_p$ 0.085\,s, B$_c$ 0.0225\,s. The step counts above are at B's default Courant
number 1.5. B$_c$ remains stable at Courant 4 (210 steps, 4.9\,s, unchanged error).
At its stability limit, B$_c$ with pure Laguerre is therefore about $5\times$
faster than A for the same basis.

\paragraph{Pitfall of the notes' projection rule (A).}
With $Q=p_z+1$ points for the projection of initial and boundary data, the error
stalls at about $2\cdot10^{-3}$ under both $p_z$- and $h$-refinement:
$3.6\cdot10^{-3}$ at $p_z=48$, and $1.8\to2.1\cdot10^{-3}$ for 20$\to$80 cells at
$p_z=64$. The nodes in $y=\mu z/2$ spread over twice the $z$-range of the basis and
under-sample the data. With $p_z+41$ points the convergence is spectral in $p_z$
($1.4\cdot10^{-2}\to2.6\cdot10^{-5}$ for $p_z=12\to48$). In $h$ the ratios are 9.9 and
18.8, close to the expected 16, until the $p_z=64$ floor of about $10^{-5}$ is
reached.

\paragraph{Long-time damped vortex, A vs.\ B.}
Density perturbation norms, Laguerre $p_z=48$, $\mu=10$.
\begin{center}
\resizebox{\textwidth}{!}{\begin{tabular}{@{}lccc@{}}
\toprule
& A & B$_c$ & B$_c$ hybrid ($N_z=24$)\\
\midrule
$\|\rho-\rho_\infty\|$, damped, $t=3$ & 0.1745 & 0.1744 & 0.1744\\
$\|\rho-\rho_\infty\|$, damped, $t=6$ & 0.03187 & 0.03188 & 0.03187\\
deviation in $z<3$, damped, $t=6$ & $2.626\cdot10^{-2}$ & $2.624\cdot10^{-2}$ & $2.625\cdot10^{-2}$\\
deviation in $z<3$, undamped, $t=3$ / $t=6$ & $7.0\cdot10^{-4}$ / $3.3\cdot10^{-3}$ & $7.3\cdot10^{-4}$ / $4.3\cdot10^{-3}$ & $2.4\cdot10^{-5}$ / $2.6\cdot10^{-5}$\\
time, undamped $+$ damped & 185\,s $+$ 168\,s & 95\,s $+$ 83\,s & 15\,s $+$ 13\,s\\
\bottomrule
\end{tabular}}
\end{center}
The two independent implementations agree to 3--4 digits on the damped dynamics,
which cross-validates the source-term treatments. The damping layer reduces the
density perturbation by a factor of 20 by $t=6$. It leaves a weak reflection of about
$2.6\cdot10^{-2}$ below the layer, which travels upstream slowly
($v-c\approx-0.18$). Without damping, the pure Laguerre bases lose accuracy as the
vortex rises: the deviation in $z<3$ grows to about $4\cdot10^{-3}$, from waves
radiated from the unresolved region. The hybrid basis keeps it at $2.6\cdot10^{-5}$.

\subsection{Conclusions on the $z$-treatment}
\begin{itemize}
\item Exact $z$-integration (Proposition 2) is mathematically clean and was verified
  to machine precision. At practical resolutions, however, it gives no accuracy
  advantage over over-integrated pointwise quadrature with the same variables and
  basis. Its costs are about $5\times$ per step, specialization to polynomial fluxes
  and $z$-only reaction, and about 50\% more code.
\item The choice of $z$-\emph{basis} dominates. A hybrid DG--Legendre $+$
  Laguerre-tail basis gives uniform resolution where the dynamics happen and is more
  accurate and much faster than any pure Laguerre configuration tested. A pure
  Laguerre basis needs many modes at large $\mu$ to resolve structures at moderate
  heights, and its resolution degrades with $z$.
\item Both codes were stable in every test, including long, under-resolved runs; no
  aliasing-driven instability was observed in B.
\item \textbf{Recommendation}: \code{laguerre\_systems\_2.cc} with the hybrid basis
  and conservative variables. Keep \code{laguerre\_systems.cc} as a reference
  implementation of the notes and as a tool to measure quadrature errors.
\end{itemize}

\paragraph{Caveats.}
All comparisons use one smooth problem family (isentropic vortices), $d=2$, a single
core, and fixed $x$-resolution (80 $Q_3$ cells). Shocks, 3D runs, other
far-field states and the role of $\mu$ were not studied. The Courant numbers of A and
B are normalized differently and are not directly comparable; step counts are.

\appendix
\section{Reproduction}
{\raggedright
\begin{itemize}
\item \code{laguerre\_systems [verify|convergence|damping|all] [n\_extra]}: logs in
  \code{results\_systems/}.
\item \code{laguerre\_systems\_2 [manufactured|cfl|vortex|damping|all]
  [conservative|primitive]}: logs in \code{results\_systems\_2/}, primitive runs with
  suffix \code{\_primitive}.
\item The density-only errors of B, the refined-quadrature runs, the
  best-approximation table and the tail mass-matrix check used small patched copies
  and scripts. These are not part of the repository; the patches changed only the
  error metric, the number of $z$-points and the driver.
\end{itemize}

\end{document}
